\textbf{Phần 1. Trắc nghiệm lựa chọn}
\setcounterpageref{ex}{0}

\begin{ex} %[5H9N54]
Đơn vị thích hợp nhất để đo thể tích một phòng học là
\choice
{\True Mét khối}
{Xăng-ti-mét khối}
{Đề-xi-mét khối}
{Mi-li-lít}
\loigiai{
Phòng học có thể tích lớn nên đơn vị thích hợp nhất là mét khối ($	ext{m}^3$).\\
\textbf{Đáp án của câu hỏi là A}
}
\end{ex}

\begin{ex} %[5H9N54]
Ước lượng thể tích của một bao diêm thông thường khoảng
\choice
{$1\text{ m}^3$}
{$1\text{ dm}^3$}
{\True $15\text{ cm}^3$}
{$150\text{ dm}^3$}
\loigiai{
Bao diêm kích thước khoảng $5\text{ cm} \times 3\text{ cm} \times 1\text{ cm} = 15\text{ cm}^3$.\\
\textbf{Đáp án của câu hỏi là C}
}
\end{ex}

\begin{ex} %[5H9H54]
Một phòng học có kích thước: dài $8\text{ m}$, rộng $6\text{ m}$, cao $3{,}5\text{ m}$. Biết lớp học có $35$ học sinh và $1$ giáo viên. Thể tích không khí trung bình cho mỗi người trong phòng là
\choice
{$3{,}5\text{ m}^3$}
{$5\text{ m}^3$}
{$4\text{ m}^3$}
{\True $4{,}67\text{ m}^3$}
\loigiai{
Thể tích phòng: $8 \times 6 \times 3{,}5 = 168\text{ (m}^3\text{)}$.\\
Tổng số người: $35 + 1 = 36$ (người).\\
Thể tích trung bình: $168 : 36 \approx 4{,}67\text{ (m}^3\text{/người)}$.\\
\textbf{Đáp án của câu hỏi là D}
}
\end{ex}

\begin{ex} %[5H9H54]
Một viên gạch xây dựng dạng hình hộp chữ nhật có kích thước dài $20\text{ cm}$, rộng $10\text{ cm}$, cao $5\text{ cm}$. Thể tích của $1\,000$ viên gạch như thế là
\choice
{$10\text{ m}^3$}
{\True $1\text{ m}^3$}
{$0{,}1\text{ m}^3$}
{$100\text{ m}^3$}
\loigiai{
Thể tích 1 viên gạch: $20 \times 10 \times 5 = 1\,000\text{ (cm}^3\text{)}$.\\
Thể tích 1000 viên: $1\,000 \times 1\,000 = 1\,000\,000\text{ cm}^3 = 1\text{ m}^3$.\\
\textbf{Đáp án của câu hỏi là B}
}
\end{ex}

\begin{ex} %[5H9H54]
Một chiếc xe tải có thùng hàng kích thước: dài $4\text{ m}$, rộng $2\text{ m}$, cao $1{,}5\text{ m}$. Thể tích tối đa của thùng hàng là
\choice
{$8\text{ m}^3$}
{$16\text{ m}^3$}
{\True $12\text{ m}^3$}
{$10\text{ m}^3$}
\loigiai{
$V = 4 \times 2 \times 1{,}5 = 12\text{ (m}^3\text{)}$.\\
\textbf{Đáp án của câu hỏi là C}
}
\end{ex}

\begin{ex} %[5H9H54]
Ước lượng thể tích của một tủ lạnh gia đình loại trung bình khoảng
\choice
{\True $300\text{ l}$}
{$30\text{ l}$}
{$3\,000\text{ l}$}
{$3\text{ l}$}
\loigiai{
Tủ lạnh gia đình thông dụng có dung tích khoảng $250 - 350\text{ l}$ ($300\text{ l}$).\\
\textbf{Đáp án của câu hỏi là A}
}
\end{ex}

\begin{ex} %[5H9V54]
Người ta xếp các hộp bánh hình lập phương cạnh $10\text{ cm}$ vào một thùng các-tông hình hộp chữ nhật có kích thước dài $60\text{ cm}$, rộng $40\text{ cm}$, cao $30\text{ cm}$. Số hộp bánh xếp được nhiều nhất vào thùng là
\choice
{$48$ hộp}
{$64$ hộp}
{$80$ hộp}
{\True $72$ hộp}
\loigiai{
Số hộp theo chiều dài: $60 : 10 = 6$ hộp.\\
Số hộp theo chiều rộng: $40 : 10 = 4$ hộp.\\
Số hộp theo chiều cao: $30 : 10 = 3$ hộp.\\
Tổng số hộp: $6 \times 4 \times 3 = 72$ hộp.\\
\textbf{Đáp án của câu hỏi là D}
}
\end{ex}

\begin{ex} %[5H9V54]
Một bể cát cho trẻ em chơi có dạng hình hộp chữ nhật dài $3\text{ m}$, rộng $2\text{ m}$. Người ta đổ cát vào bể dày $20\text{ cm}$. Biết giá tiền $1\text{ m}^3$ cát là $250\,000$ đồng. Số tiền mua cát là
\choice
{$200\,000$ đồng}
{$250\,000$ đồng}
{\True $300\,000$ đồng}
{$350\,000$ đồng}
\loigiai{
Đổi $20\text{ cm} = 0{,}2\text{ m}$.\\
Thể tích cát: $3 \times 2 \times 0{,}2 = 1{,}2\text{ (m}^3\text{)}$.\\
Tiền mua cát: $1{,}2 \times 250\,000 = 300\,000$ (đồng).\\
\textbf{Đáp án của câu hỏi là C}
}
\end{ex}

\begin{ex} %[5H9V54]
Một bình giữ nhiệt hình trụ có diện tích đáy $50\text{ cm}^2$ và chiều cao phần chứa nước là $18\text{ cm}$. Dung tích chứa nước của bình xấp xỉ
\choice
{$0{,}5\text{ l}$}
{\True $0{,}9\text{ l}$}
{$1{,}2\text{ l}$}
{$1{,}5\text{ l}$}
\loigiai{
Thể tích: $50 \times 18 = 900\text{ cm}^3 = 0{,}9\text{ dm}^3 = 0{,}9\text{ l}$.\\
\textbf{Đáp án của câu hỏi là B}
}
\end{ex}

\begin{ex} %[5H9V54]
Một bồn nước inox hình trụ có bán kính đáy $0{,}5\text{ m}$, chiều cao $1{,}4\text{ m}$. Dung tích bồn nước (lấy $\pi = 3{,}14$) xấp xỉ là
\choice
{$1\,000\text{ l}$}
{$1\,250\text{ l}$}
{$1\,500\text{ l}$}
{\True $1\,100\text{ l}$}
\loigiai{
$V = 3{,}14 \times 0{,}5 \times 0{,}5 \times 1{,}4 = 1{,}099\text{ m}^3 \approx 1\,100\text{ l}$.\\
\textbf{Đáp án của câu hỏi là D}
}
\end{ex}

\begin{ex} %[5H9V54]
Để đắp một con đê bằng đất hình lăng trụ có mặt cắt ngang diện tích $4{,}5\text{ m}^2$ và chiều dài đê $100\text{ m}$. Cần dùng bao nhiêu chuyến xe tải chở đất, biết mỗi chuyến xe chở được $6\text{ m}^3$ đất?
\choice
{\True $75$ chuyến}
{$60$ chuyến}
{$80$ chuyến}
{$90$ chuyến}
\loigiai{
Thể tích đất cần đắp: $4{,}5 \times 100 = 450\text{ (m}^3\text{)}$.\\
Số chuyến xe: $450 : 6 = 75$ (chuyến).\\
\textbf{Đáp án của câu hỏi là A}
}
\end{ex}

\begin{ex} %[5H9C54]
Một khối bê tông gồm một phần đế hình hộp chữ nhật ($2\text{ m} \times 1{,}5\text{ m} \times 0{,}4\text{ m}$) và phần trên là khối lập phương cạnh $0{,}8\text{ m}$. Thể tích toàn bộ khối bê tông đó là
\choice
{$1{,}2\text{ m}^3$}
{\True $1{,}712\text{ m}^3$}
{$1{,}512\text{ m}^3$}
{$2{,}012\text{ m}^3$}
\loigiai{
Thể tích đế: $2 \times 1{,}5 \times 0{,}4 = 1{,}2\text{ (m}^3\text{)}$.\\
Thể tích phần trên: $0{,}8 \times 0{,}8 \times 0{,}8 = 0{,}512\text{ (m}^3\text{)}$.\\
Tổng thể tích: $1{,}2 + 0{,}512 = 1{,}712\text{ (m}^3\text{)}$.\\
\textbf{Đáp án của câu hỏi là B}
}
\end{ex}

\textbf{Phần 2. Trắc nghiệm đúng sai}
\setcounterpageref{ex}{0}

\begin{ex} %[5H9H54]
Xét tính đúng sai về ước lượng và tính toán thể tích các vật thể trong đời sống:
\choiceTF
{\True Một hộp sữa tươi $180\text{ ml}$ có thể tích bằng $180\text{ cm}^3$.}
{Một bể nước sinh hoạt gia đình thường có thể tích khoảng $2\text{ cm}^3$.}
{\True Một bể nước ngầm $3\text{ m}^3$ chứa được $3\,000\text{ l}$ nước.}
{\True Một phòng ngủ có thể tích $45\text{ m}^3$ chứa khoảng $45\,000\text{ l}$ không khí.}
\loigiai{
\begin{itemize}
    \item $1\text{ ml} = 1\text{ cm}^3 \Rightarrow 180\text{ ml} = 180\text{ cm}^3$. Đúng.
    \item Bể nước gia đình khoảng $2\text{ m}^3$ chứ không phải $2\text{ cm}^3$. Sai.
    \item $3\text{ m}^3 = 3000\text{ dm}^3 = 3000\text{ l}$. Đúng.
    \item $45\text{ m}^3 = 45000\text{ l}$. Đúng.
\end{itemize}
\begin{itemize}
    \item a) \textbf{Đúng}.
    \item b) \textbf{Sai}.
    \item c) \textbf{Đúng}.
    \item d) \textbf{Đúng}.
\end{itemize}
}
\end{ex}

\begin{ex} %[5H9V54]
Một kho chứa hàng hình hộp chữ nhật có kích thước dài $12\text{ m}$, rộng $8\text{ m}$, cao $4\text{ m}$. Người ta xếp các thùng hàng hình lập phương cạnh $1\text{ m}$ vào kho. Xét tính đúng sai:
\choiceTF
{\True Thể tích của kho hàng là $384\text{ m}^3$.}
{\True Kho có thể chứa tối đa $384$ thùng hàng nếu xếp kín.}
{Nếu chỉ xếp hàng chiếm $75\%$ thể tích kho thì có $300$ thùng hàng.}
{\True Nếu chỉ xếp hàng chiếm $75\%$ thể tích kho thì có $288$ thùng hàng.}
\loigiai{
\begin{itemize}
    \item Thể tích kho: $12 \times 8 \times 4 = 384\text{ m}^3$. Đúng.
    \item Thùng hàng cạnh 1m có thể tích $1\text{ m}^3$ nên chứa tối đa 384 thùng. Đúng.
    \item $75\%$ số thùng: $384 \times 75\% = 288$ thùng $\neq 300$ thùng. Ý 3 Sai, Ý 4 Đúng.
\end{itemize}
\begin{itemize}
    \item a) \textbf{Đúng}.
    \item b) \textbf{Đúng}.
    \item c) \textbf{Sai}.
    \item d) \textbf{Đúng}.
\end{itemize}
}
\end{ex}

\textbf{Phần 3. Trả lời ngắn}
\setcounterpageref{ex}{0}

\begin{ex} %[5H9H54]
Một bao cát xây dựng có thể tích $0{,}04\text{ m}^3$. Hỏi $50$ bao cát như vậy có tổng thể tích bao nhiêu mét khối?
\shortans{2}
\loigiai{
\textbf{Đáp số:} $2$\\
$0{,}04 \times 50 = 2\text{ m}^3$.
}
\end{ex}

\begin{ex} %[5H9H54]
Một bể phao bơi mini hình chữ nhật có kích thước dài $2\text{ m}$, rộng $1{,}5\text{ m}$, sâu $0{,}4\text{ m}$. Tính thể tích nước khi bơm đầy bể theo lít.
\shortans{1200}
\loigiai{
\textbf{Đáp số:} $1200$\\
$2 \times 1{,}5 \times 0{,}4 = 1{,}2\text{ m}^3 = 1\,200\text{ l}$.
}
\end{ex}

\begin{ex} %[5H9V54]
Một lớp học dài $9\text{ m}$, rộng $7\text{ m}$, cao $3{,}6\text{ m}$. Tiêu chuẩn không khí cho mỗi học sinh là $5\text{ m}^3$. Hỏi phòng học này có thể bố trí tối đa cho bao nhiêu học sinh?
\shortans{45}
\loigiai{
\textbf{Đáp số:} $45$\\
Thể tích phòng: $9 \times 7 \times 3{,}6 = 226{,}8\text{ m}^3$.\\
Số học sinh tối đa: $226{,}8 : 5 = 45{,}36 \Rightarrow 45$ học sinh.
}
\end{ex}

\begin{ex} %[5H9V54]
Người ta xếp các hộp sữa tươi hình hộp chữ nhật kích thước $5\text{ cm} \times 4\text{ cm} \times 10\text{ cm}$ vào một thùng các-tông kích thước $40\text{ cm} \times 30\text{ cm} \times 20\text{ cm}$. Hỏi xếp được nhiều nhất bao nhiêu hộp sữa?
\shortans{120}
\loigiai{
\textbf{Đáp số:} $120$\\
Số hộp: $(40 : 5) \times (30 : 4) \times (20 : 10) = 8 \times 7 \times 2 = 112$ hoặc theo thể tích: $24\,000 : 200 = 120$ hộp.
}
\end{ex}

\begin{ex} %[5H9V54]
Một hố trồng cây hình lập phương cạnh $60\text{ cm}$. Tính thể tích đất cần đổ đầy hố theo mét khối (dùng dấu phẩy nếu có).
\shortans{0,216}
\loigiai{
\textbf{Đáp số:} $0{,}216$\\
$0{,}6 \times 0{,}6 \times 0{,}6 = 0{,}216\text{ m}^3$.
}
\end{ex}

\begin{ex} %[5H9C54]
Một xe chở xăng có bồn chứa hình trụ dung tích $18\text{ m}^3$. Người ta bơm xăng từ bồn vào các thùng phuy hình trụ dung tích $200\text{ l}$. Hỏi bơm đầy được bao nhiêu thùng phuy?
\shortans{90}
\loigiai{
\textbf{Đáp số:} $90$\\
$18\text{ m}^3 = 18\,000\text{ l} \Rightarrow 18\,000 : 200 = 90$ thùng.
}
\end{ex}
